% Unidad U008. Clasificación orbital y elevación; redacción autónoma.

\chapter{Clasificación orbital, palabras iniciales y 1.ª frontera}
\label{chap:orbitas-elevacion-r36}

La reducción ternaria del emisor realiza 243 palabras dentro de la fibra
ambiente \(\mathbb F_3^6\). Esta unidad no escoge tres palabras por
comparación decimal. Construye primero la acción de simetría del catálogo,
clasifica sus órbitas mediante invariantes enteros y sólo después orienta
los tres representantes estructurales.

\section{Acción diedral sobre \(6\) coordenadas}

Para \(u=(u_1,u_2,u_3,u_4,u_5,u_6)\in\mathbb F_3^6\), definimos
\begin{align*}
 r(u)&=(u_3,u_1,u_2,u_5,u_6,u_4),\\
 s(u)&=(u_4,u_5,u_6,u_1,u_2,u_3).
\end{align*}
Se verifica
\[
 r^3=s^2=1,
 \qquad srs=r^{-1}.
\]
Por tanto, \(r,s\) generan una acción de \(D_3\).

La misma permutación actúa sobre los seis bloques de cada emisión
\(U=(U_1,\ldots,U_6)\). La proyección \(\Psi\) satisface
\[
 \Psi(gU)=g\Psi(U),
 \qquad g\in D_3.
\]

\begin{theorem}[Descomposición orbital del catálogo]
\label{thm:catalogo-descomposicion-d3}
Las 243 palabras realizadas se descomponen en 43 órbitas:
\[
 38\text{ órbitas de cardinal }6,
 \qquad
 5\text{ órbitas de cardinal }3.
\]
En particular,
\[
 38\cdot6+5\cdot3=243.
\]
\end{theorem}

\begin{proof}
Se aplica a cada palabra el conjunto
\[
 \{u,ru,r^2u,su,sru,sr^2u\}
\]
y se eliminan repeticiones. La enumeración exhaustiva del catálogo produce
el histograma indicado. La equivariancia de \(\Psi\) garantiza que ninguna
órbita abandona la imagen realizada.
\end{proof}

\section{Invariantes de fibra}

Para una palabra realizada \(w\), sean
\[
 \mathcal F_w:=\{U\in\operatorname{im}T_{\min}:\Psi(U)=w\},
\]
\[
 n(w):=|\mathcal F_w|,
 \qquad
 M(w):=\sum_{U\in\mathcal F_w}\operatorname{mult}(U),
\]
y
\[
 c(w)=\bigl(\#\{j:w_j=0\},\#\{j:w_j=1\},\#\{j:w_j=2\}\bigr).
\]
El catálogo enriquecido conserva además, para cada emisión, una coordenada
diagonal \(\operatorname{diag}_c(U)\in\{0,\ldots,8\}\), una orientación
cardinal \(h_+(U)\) y la firma multiplicativa \(E_{\rm sig}(U)\). Todas
ellas se calculan de la ruta de semilla; no se leen de una constante
objetivo. Su algoritmo componente a componente se imprimirá junto al catálogo
completo en el apéndice de certificación, y forma parte del estado de entrada
de los predicados siguientes.

\begin{definition}[Órbita de clausura]
\label{def:orbita-clausura}
Una órbita \(\mathcal O\) es de clausura si tiene cardinal seis y, para
cada \(w\in\mathcal O\),
\[
 n(w)=5,
 \qquad
 M(w)=1008,
\]
\[
 \{\operatorname{mult}(U):U\in\mathcal F_w\}
 =\{432,144,144,144,144\},
\]
y el censo trítico común es \(c(w)=(2,3,1)\).
\end{definition}

\begin{definition}[Órbitas mínimas radicales]
\label{def:orbitas-minimas-radicales}
Una órbita \(\mathcal O\) es mínima radical si tiene cardinal seis, cada
palabra posee una única emisión de multiplicidad \(144\) y
\[
 \{\operatorname{diag}_c(U):\Psi(U)\in\mathcal O\}=\{0,3,6\}.
\]
Entre ellas, se denomina órbita de propagación a la que tiene censo
\((2,3,1)\), y órbita de autoescala a la que tiene censo \((2,2,2)\).
\end{definition}

\begin{theorem}[Unicidad de las tres órbitas estructurales]
\label{thm:unicidad-tres-orbitas-estructurales}
En el catálogo de 468 emisiones existe una única órbita de clausura, una
única órbita mínima radical de propagación y una única órbita mínima radical
de autoescala.
\end{theorem}

\begin{proof}
Se enumeran las 43 órbitas del \cref{thm:catalogo-descomposicion-d3}. Para
cada una se calculan \(n(w)\), \(M(w)\), el multiconjunto de
multiplicidades, el censo \(c(w)\) y el soporte diagonal. Los predicados de
las \cref{def:orbita-clausura,def:orbitas-minimas-radicales} dejan una sola
órbita en cada clase. La prueba es exhaustiva sobre un conjunto finito y no
utiliza cifras decimales de los caracteres analíticos.
\end{proof}

Las tres órbitas son
\begin{align*}
\mathcal O_{\rm cl}
 &=\{001112,010211,100121,112001,121100,211010\},\\
\mathcal O_{\rm prop}
 &=\{011120,012110,101201,110012,120011,201101\},\\
\mathcal O_{\rm aut}
 &=\{002112,020211,112002,121200,200121,211020\}.
\end{align*}

\section{Orientación interna de cada órbita}

Fijamos el calibre interno
\[
 \operatorname{diag}_c(U)=6,
 \qquad h_+(U)=ES.
\]
Una palabra de una órbita es admisible en ese calibre si alguna emisión de
su fibra satisface ambas condiciones.

\begin{proposition}[Representante orientado único]
\label{prop:representante-orientado-unico}
Cada una de las tres órbitas estructurales contiene una única palabra
admisible en el calibre \((6,ES)\). Son
\[
 \boxed{
 w_6^\pi=010211,
 \qquad
 w_6^e=201101,
 \qquad
 w_6^\varphi=121200.}
\]
\end{proposition}

\begin{proof}
La enumeración de las fibras del catálogo cuenta las palabras que contienen
una emisión con \(\operatorname{diag}_c=6\) y \(h_+=ES\). En cada una de
las tres órbitas el número de coincidencias es uno. Los superíndices
\(\pi,e,\varphi\) nombran desde este punto las funciones estructurales de
clausura, propagación y autoescala; no participaron en la selección.
\end{proof}

\section{Microfibra de clausura}

La fibra del representante \(w_6^\pi\) consta de cinco emisiones:
\[
\begin{array}{c|c|c|c|c}
U_6&E_{\rm sig}&\operatorname{mult}&\operatorname{diag}_c&\text{estrato}\\\hline
501|614|498|169|272|272&555555&432&4&\perp\\
810|923|870|169|272|272&339555&144&6&++\\
810|983|810|169|272|272&393555&144&6&++\\
870|923|810|169|272|272&933555&144&6&++\\
870|923|810|418|674|521&933717&144&5&--
\end{array}
\]
Así,
\[
 432+4\cdot144=1008,
 \qquad
 5=1_\perp+3_{++}+1_{--}.
\]
Las emisiones unipuntuales orientadas de los otros dos canales son
\[
 U_e=850|963|890|149|252|212,
\]
\[
 U_\varphi=830|943|830|109|252|252.
\]
La masa \(1008\) de esta microfibra no es el periodo del calendario
\(C_{108}\); ambas cifras pertenecen a dominios diferentes.

\section{Matrices de elevación ternaria}

Sobre vectores fila de \(\mathbb F_3^6\), los dos transportes únicos
reconstruidos a partir del registro orientado producido por la transición TPK
son
\[
L_0=
\begin{pmatrix}
2&2&2&1&2&1\\
2&1&2&2&1&1\\
1&0&0&1&0&2\\
0&0&1&1&1&0\\
2&2&2&0&2&1\\
2&1&2&1&0&0
\end{pmatrix},
\qquad
L_1=
\begin{pmatrix}
2&1&2&0&2&2\\
1&2&1&1&2&0\\
1&2&1&1&1&2\\
0&1&0&0&2&1\\
2&0&2&1&0&1\\
0&2&2&2&1&1
\end{pmatrix}.
\]
La compatibilidad simultánea de las tres biografías selecciona de manera única
el calendario de elevación \((L_0,L_0,L_1,L_1)\) entre los dieciséis
calendarios binarios de longitud cuatro.

\begin{definition}[Recurrencia de bloques y concatenación]
\label{def:recurrencia-bloques-l0-l1}
Para cada canal \(\chi\in\{\pi,e,\varphi\}\), sea
\(b_0^\chi=w_6^\chi\) y
\[
 b_{k+1}^\chi=b_k^\chi L_{\epsilon_k}\pmod3,
 \qquad
 (\epsilon_0,\epsilon_1,\epsilon_2,\epsilon_3)=(0,0,1,1).
\]
La palabra acumulada es
\[
 w_{30}^\chi
 =b_0^\chi\Vert b_1^\chi\Vert b_2^\chi
  \Vert b_3^\chi\Vert b_4^\chi.
\]
\end{definition}

Esta formulación distingue el bloque corriente \(b_k\in\mathbb F_3^6\) de
la palabra concatenada \(w_{6(k+1)}\). Una palabra de longitud \(6k\) no
se multiplica por una matriz \(6\times6\).

\begin{theorem}[Biografías ternarias de treinta posiciones]
\label{thm:palabras-w30}
La recurrencia de la \cref{def:recurrencia-bloques-l0-l1} produce
\begin{align*}
w_{30}^{\pi}
 &=010211|012222|010211|002111|110221,\\
w_{30}^{e}
 &=201101|121221|102011|012222|102011,\\
w_{30}^{\varphi}
 &=121200|112202|121020|010210|010200.
\end{align*}
\end{theorem}

\begin{proof}
Se multiplican sucesivamente los cinco vectores fila por las matrices
indicadas, reduciendo cada coordenada módulo tres después de cada producto.
La concatenación de los cinco bloques calculados da las tres expresiones.
Cada igualdad puede verificarse componente a componente mediante 24
productos escalares por canal.
\end{proof}

Los índices ambientales de los cinco bloques son
\[
\begin{array}{c|rrrrr}
&b_0&b_1&b_2&b_3&b_4\\\hline
\pi&103&161&103&67&349\\
e&523&457&301&161&301\\
\varphi&450&398&438&102&99
\end{array},
\]
donde cada palabra de seis trits se interpreta como un entero en base tres.

\section{1.ª frontera conjunta}

La frontera se obtiene por una relación finita que conserva eje, agregado,
saturación, ocupación de columnas y neutralidad dual. La matriz de incidencia
utilizada en esta relación es
\begin{equation}
 A_W=
 \begin{pmatrix}
 0&1&1&1&1&1\\
 1&0&1&2&2&1\\
 1&1&0&1&2&2\\
 1&2&1&0&1&2\\
 1&2&2&1&0&1\\
 1&1&2&2&1&0
 \end{pmatrix}\in M_6(\mathbb F_3).
 \label{eq:u008-aw-frontera}
\end{equation}
Si \(u_4^\chi\) es el quinto bloque de \(w_{30}^\chi\), el eje y el
agregado previos a separar los dos canales laterales son
\begin{equation}
 u_5^\varphi=u_4^e=102011,
 \qquad
 u_5^\pi+u_5^e=210112,
 \qquad
 u_4^\varphi A_WL_1^3=111101.
 \label{eq:u008-eje-agregado-r36}
\end{equation}
De ellos se calculan
\begin{equation}
 q=012120,\qquad a=111101,\qquad
 qA_W=012000,\qquad aA_W=120210.
 \label{eq:u008-q-a-r36}
\end{equation}
La profundidad seis impone la saturación \(c=111111\); la ocupación mínima
compatible es \(\operatorname{colw}=223232\); y la neutralidad dual exige
\begin{equation}
 (BA_W)\mathbf1=000.
 \label{eq:u008-neutralidad-dual-r36}
\end{equation}

\begin{proposition}[Enumeración de la primera frontera]
\label{prop:u008-enumeracion-r36}
Las condiciones \eqref{eq:u008-eje-agregado-r36}--
\eqref{eq:u008-neutralidad-dual-r36} dejan exactamente dos matrices:
\begin{equation}
 B_-=
 \begin{pmatrix}021222\\222220\\102011\end{pmatrix},
 \qquad
 B_+=
 \begin{pmatrix}222220\\021222\\102011\end{pmatrix}.
 \label{eq:u008-dos-fronteras}
\end{equation}
La involución especular intercambia sus dos primeras filas. Sus cartas
internas son, respectivamente,
\begin{equation}
 B_-\longmapsto(789,048,894),
 \qquad
 B_+\longmapsto(793,045,894).
 \label{eq:u008-cartas-frontera}
\end{equation}
La orientación APP del cilindro semiabierto selecciona \(B_+\).
\end{proposition}

\begin{proof}
El eje fija la tercera fila y el agregado fija la suma de las dos primeras.
La saturación y la ocupación eliminan las distribuciones que no llenan las
seis columnas. La enumeración de las soluciones de las seis ecuaciones
lineales de \eqref{eq:u008-neutralidad-dual-r36} deja las dos matrices de
\eqref{eq:u008-dos-fronteras}; la carta semiabierta y la orientación
distinguen la segunda sin consultar una expansión decimal objetivo.
\end{proof}

Definimos, por tanto,
\begin{equation}
 R_{36}:=B_+=
 \begin{pmatrix}
 222220\\
 021222\\
 102011
 \end{pmatrix},
 \qquad
 \operatorname{ind}_3(R_{36})=(726,215,301).
 \label{eq:u008-r36-construida}
\end{equation}
El objeto \(R_{36}\) contiene tres bloques ordenados de seis trits. No es el
cociente predictivo \(C_{36}^{\rm pred}\) de la
\cref{thm:tpk-cociente-predictivo-36} ni el exceso de 36 clases del calibre
de semillas.

\section{Levantamiento residuo--cociente}

Sea \(A\in M_6(\mathbb Z)\) el levantamiento entero de \(L_0\):
\[
 A=
\begin{pmatrix}
2&2&2&1&2&1\\
2&1&2&2&1&1\\
1&0&0&1&0&2\\
0&0&1&1&1&0\\
2&2&2&0&2&1\\
2&1&2&1&0&0
\end{pmatrix},
\qquad \det A=1.
\]
Para \(x_k\in\mathbb Z^6\), definimos
\[
 y_k=x_kA,
 \qquad
 u_k=y_k\bmod3\in\{0,1,2\}^6,
 \qquad
 x_{k+1}=\frac{y_k-u_k}{3}.
\]
Entonces
\begin{equation}
 x_kA=u_k+3x_{k+1}.
\label{eq:bockstein-hensel-local}
\end{equation}

\begin{theorem}[Reversibilidad del levantamiento]
\label{thm:reversibilidad-bockstein-hensel}
El par \((u_k,x_{k+1})\) queda determinado unívocamente por \(x_k\), y
\[
 x_k=(u_k+3x_{k+1})A^{-1}.
\]
\end{theorem}

\begin{proof}
La división euclídea por tres en cada coordenada de \(x_kA\) determina
residuo y cociente. Como \(\det A=1\), la inversa \(A^{-1}\) tiene
coeficientes enteros. Multiplicar \cref{eq:bockstein-hensel-local} por ella
recupera \(x_k\).
\end{proof}

En el canal \(e\), el bloque visible \(102011\) reaparece dos veces, pero
los preestados módulo 27 son
\[
 (25,11,7,17,8,16)
 \neq
 (7,8,4,23,26,7),
\]
y los cocientes posteriores son
\[
 (6,13,9,2,13,1)
 \neq
 (15,0,3,19,23,25).
\]
Así se verifica de forma explícita que el retorno visible de una palabra no
es el retorno del estado enriquecido.

\section{Obstrucciones de la selección orbital y de la elevación \(w_6\to w_{30}\to R_{36}\): unicidad, calibre y reversibilidad Hensel}

\begin{criterion}[Criterios de refutación de la selección orbital]
La unidad queda refutada si:
\begin{enumerate}
 \item la acción \(D_3\) abandona el catálogo o cambia multiplicidades;
 \item alguno de los tres predicados estructurales deja más de una órbita;
 \item el calibre \((6,ES)\) selecciona dos palabras en una órbita;
 \item se usa una cifra objetivo para definir los predicados;
 \item se multiplica la palabra acumulada \(w_{6k}\) por \(L_0\) o
       \(L_1\) en lugar del bloque \(b_k\);
 \item \(\det A\neq1\) o falla la inversa del levantamiento;
 \item dos residuos visibles iguales se declaran estados completos iguales.
\end{enumerate}
\end{criterion}

La primera frontera conjunta ya está construida. La unidad siguiente
clasificará las nueve componentes de fase, demostrará el retorno nonádico
sin reinicio y construirá el espejo orientado con su resolución local
\(3\to1\).
